 \item 梯度分析法在本实验中表现最好，能够有效估计运动模糊的方向和长度
    \item 盲反卷积法虽然计算复杂度高，但在没有先验信息的情况下能取得较好效果
    \item 维纳滤波的参数调整对复原质量有显著影响，需要根据具体情况进行优化
    \item 单一评价指标不足以全面反映复原效果，应结合多种指标综合评价
\end{enumerate}

在实际应用中，可以根据图像的特点和模糊类型选择合适的估计方法。对于运动模糊，梯度分析法和盲反卷积法通常能取得较好效果；对于其他类型的模糊，如高斯模糊或散焦模糊，可能需要其他专门的估计方法。

维纳滤波虽然是一种经典有效的复原方法，但在强噪声环境下可能不够稳健。未来可以尝试结合深度学习的方法，如卷积神经网络(CNN)或生成对抗网络(GAN)来进行模糊图像复原，这些方法在处理复杂模糊和噪声方面可能有更好的表现。

\section{对本实验过程及方法、手段的改进建议:}
\begin{enumerate}
    \item 可以尝试更多类型的模糊图像，包括不同方向、不同长度的运动模糊，以及其他类型的模糊（如高斯模糊、散焦模糊等）
    \item 引入更多的模糊核估计方法，如基于卷积神经网络的估计方法
    \item 改进维纳滤波，如考虑非均匀噪声或空变噪声的情况
    \item 增加更多的评价指标，如感知质量评分(MOS)、无参考图像质量评估(NIQE)等
    \item 设计更高效的参数自动优化算法，如贝叶斯优化或进化算法
    \item 尝试结合多种方法的优点，如先用盲反卷积估计初始核，再用频谱分析或梯度分析进行精细调整
\end{enumerate}